\documentclass{article}

\usepackage{amsmath}

\usepackage{listings}
\usepackage{microtype}


\begin{document}

    \section{Einleitung}
    \label{sec:einleitung}

    \subsection{Motivation: Die Notwendigkeit physikalisch korrekter Raumakustik}
    \label{subsec:motivation}

    Die räumliche Wahrnehmung von Schallereignissen basiert auf der Auswertung komplexer Reflexions- und Absorptionsmuster durch das menschliche Gehör.
    In Anwendungsbereichen wie der digitalen Audioproduktion,
    der Architekturakustik sowie immersiven Umgebungen (z.\,B. Virtual Reality) ist eine physikalisch akkurate Reproduktion dieser Raumakustik maßgeblich für die Authentizität der Simulation.

    Etablierte Verfahren der digitalen Signalverarbeitung (DSP) verwenden zur Raumsimulation traditionell algorithmische Halleffekte (Reverbs).
    Diese Systeme generieren späte Reflexionen (\textit{Late Reverberation}) durch Delay-Netzwerke und Feedback-Schleifen, wodurch das zeitliche Abklingen der Schallenergie künstlich angenähert wird.
    Da diese Algorithmen nicht auf realen Geometrien basieren, können sie raumspezifische physikalische Phänomene wie die frequenz- und winkelabhängige Materialabsorption,
    den entfernungsbedingten Energieabfall (\textit{Inverse Square Law}) sowie diskrete Phasenauslöschungen (Kammfiltereffekte) nicht nativ abbilden.


    Um eine höhere physikalische Genauigkeit zu erzielen, ist die Berechnung der exakten Schallausbreitung im Raum erforderlich.
    Das akustische Raytracing simuliert hierbei die Ausbreitungswege von Schallwellen geometrisch.
    Die praktische Umsetzung dieses deterministischen und stochastischen Verfahrens scheiterte in der Vergangenheit oftmals an der hohen Rechenkomplexität,
    da für eine realistische Dichte des simulierten Schallfeldes Millionen von Strahlengängen berechnet werden müssen.
    Durch die Nutzung moderner Multi-Core-Architekturen und hochoptimierter, nebenläufiger (multithreaded) Software-Designs lassen sich diese Berechnungen mittlerweile effizient auf lokalen Systemen durchführen.
    Dies motiviert die Entwicklung von Systemen, die eine physikalisch korrekte Berechnung der Raumgeometrie mit einer direkten DSP-Anwendung auf Audiosignale verbinden.



    \subsection{Zielsetzung}
    \label{subsec:zielsetzung}

    Das Hauptziel dieser Arbeit ist die Konzeption und Implementierung einer hocheffizienten, nebenläufigen (multithreaded) Raytracing-Engine zur physikalisch korrekten Simulation von Raumakustik.
    Um eine performante und modular erweiterbare Pipeline zu schaffen, werden folgende Teilziele definiert:

    \begin{enumerate}
        \item \textbf{Physikalisch akkurate Akustiksimulation:}\\
        Im Zentrum der Entwicklung steht die deterministische und stochastische Berechnung früher Reflexionen (\textit{Early Reflections}) und der späten Hallfahne (\textit{Late Reverberation}).
        Das System muss grundlegende physikalische Gesetzmäßigkeiten nativ abbilden, darunter den entfernungsbedingten Energieabfall (\textit{Inverse Square Law}),
        material- und winkelabhängige Absorptionskoeffizienten sowie laufzeitbedingte Signalverschmierungen (\textit{Temporal Smearing}).
        Die Ausgabe dieser Simulation erfolgt in Form einer diskreten Raumimpulsantwort (Impulse Response, IR).

        \item \textbf{Maximierung der Rechenleistung durch CPU-Multithreading:}\\
        Ein zentrales technisches Ziel ist die Überwindung bisheriger Performance-Engpässe bei der Berechnung von Millionen von Schallstrahlen.
        Anstelle eines netzwerkbasierten, verteilten Systems wird ein lokaler, hochgradig parallelisierter Ansatz verfolgt.
        Durch den Einsatz der Programmiersprache \texttt{Rust} und der Datenparallelisierungs-Bibliothek \texttt{rayon} soll die CPU-Auslastung durch dynamisches \textit{Work-Stealing} maximiert werden.
        Dies gewährleistet eine speichersichere und recheneffiziente Ausführung der Monte-Carlo-Integration ohne den Overhead klassischer Cluster-Frameworks.

        \item \textbf{Strikte architektonische Entkopplung:}\\
        Die Systemarchitektur verlangt eine klare Trennung der Domänen Physik und Audio.
        Die rechenintensive Berechnung der Strahlengänge und Raumgeometrien wird vollständig in Rust gekapselt.
        Die Orchestrierung der Pipeline, die Aggregation der JSON-basierten Simulationsdaten, die anschließende digitale Signalverarbeitung (DSP) sowie die rechenintensive Faltung
        (\textit{Convolution}) der berechneten Impulsantworten mit trockenen Audiosignalen erfolgen in Python.
        Diese strikte Entkopplung stellt sicher, dass das Geometrie-Backend (2D/3D) unabhängig von der Audioverarbeitung weiterentwickelt werden kann.
    \end{enumerate}

    \subsection{Aufbau der Arbeit}
    \label{subsec:aufbau_der_arbeit}

    Um die genannten Zielsetzungen systematisch zu adressieren, ist die vorliegende Arbeit wie folgt strukturiert:

    \paragraph{Kapitel 2} legt das theoretische Fundament.
    Zunächst werden die akustischen Prinzipien der Schallausbreitung im Raum erläutert, darunter das Abstandsgesetz (\textit{Inverse Square Law}) und die Materialabsorption.
    Darauf aufbauend wird das stochastische Verfahren der Monte-Carlo-Integration im Kontext des Raytracings beschrieben.
    Den Abschluss des Kapitels bilden die Grundlagen der digitalen Signalverarbeitung (DSP), insbesondere die Struktur von Raumimpulsantworten und deren Anwendung durch schnelle Faltung
    (\textit{Fast Fourier Transform}).

    \paragraph{Kapitel 3} widmet sich der Konzeption der Systemarchitektur und der Orchestrierung der gesamten Pipeline.
    Es begründet die technologische und architektonische Entkopplung der rechenintensiven Physiksimulation von der Audiosignalverarbeitung und detailliert den prozessübergreifenden Datenfluss.

    \paragraph{Kapitel 4} beschreibt die spezifische Implementierung der akustischen Physik-Engine in Rust.
    Hierzu zählen die Vektormathematik der 2D-Schnittpunkttests, die Datenparallelisierung mittels \texttt{rayon} sowie Architekturmuster zur Vermeidung von Speicherallokations-Engpässen.

    \paragraph{Kapitel 5} behandelt die Orchestrierung und digitale Signalverarbeitung in Python.
    Es wird erläutert, wie die rohen Simulationsdaten eingelesen und mittels der \texttt{SciPy}-Bibliothek verarbeitet werden,
    um die finale Faltung der berechneten Impulsantworten mit trockenen Audiosignalen durchzuführen.

    \paragraph{Kapitel 6} evaluiert und verifiziert das entwickelte System.
    Die physikalische Korrektheit der Simulation wird durch visuelles Debugging der Strahlengänge (\textit{Ray Path Plotting}) sowie durch die Analyse der finalen Audiowellenformen validiert.

    \paragraph{Kapitel 7} fasst die Ergebnisse zusammen, reflektiert die Erreichung der Zielsetzungen und gibt einen Ausblick auf zukünftige Erweiterungsmöglichkeiten,
    wie etwa die Migration in eine vollständige 3D-Volumengeometrie.
\end{document}
